%% 
%% Copyright 2019-2020 Elsevier Ltd
%% This file is part of the 'CAS Bundle'.
%% --------------------------------------
%% 
%% It may be distributed under the conditions of the LaTeX Project Public
%% License, either version 1.2 of this license or (at your option) any
%% later version.  The latest version of this license is in
%%    http://www.latex-project.org/lppl.txt
%% and version 1.2 or later is part of all distributions of LaTeX
%% version 1999/12/01 or later.
%% The list of all files belonging to the 'CAS Bundle' is
%% given in the file `manifest.txt'. 
%% Template article for cas-sc documentclass for 
%% double column output.
%\PassOptionsToPackage{numbers,sort&compress}{natbib}

\documentclass[a4paper,fleqn]{cas-sc}

%\documentclass[a4paper,fleqn,authoryear]{cas-sc}
%\documentclass[a4paper,fleqn,longmktitle]{cas-sc}

\usepackage[square,numbers,sort&compress]{natbib}
\bibliographystyle{model1-num-names}

\usepackage{placeins}
\usepackage{array}
\usepackage{makecell}
\usepackage{subcaption}
\usepackage[range-units=single, per-mode=symbol]{siunitx}
\sisetup{
  group-separator = {,},
  group-minimum-digits = 4
}
% \sisetup{
%   range-phrase = -- % Sets the separator to an en-dash
% }

\usepackage[printonlyused]{acronym}
\usepackage{tabularx}
\usepackage{booktabs}
% \usepackage{biblatex}
\usepackage{enumitem}

\usepackage[en-GB, showseconds=false]{datetime2}
\DTMsetstyle{iso}
% \usepackage{appendix}
% \usepackage{float}
% \usepackage{stfloats}
% \usepackage{dblfloatfix}

\newcommand\mycite[1]{\citeauthor{#1}~(\citeyear{#1})}

%%% Handy commands for image/figure insertion
\newcommand*{\RootPicDir}{pic}
\newcommand*{\PicDir}{\RootPicDir}
\newcommand*{\ResetPicDir}{\renewcommand*{\PicDir}{\RootPicDir}}
\newcommand*{\SetPicSubDir}[1]{\renewcommand*{\PicDir}{\RootPicDir /#1}}
\newcommand*{\Pic}[2]{\PicDir /#2.#1}

%%%Author definitions
\def\tsc#1{\csdef{#1}{\textsc{\lowercase{#1}}\xspace}}
\tsc{WGM}
\tsc{QE}
\tsc{EP}
\tsc{PMS}
\tsc{BEC}
\tsc{DE}
%%%

% Uncomment and use as if needed
%\newtheorem{theorem}{Theorem}
%\newtheorem{lemma}[theorem]{Lemma}
%\newdefinition{rmk}{Remark}
%\newproof{pf}{Proof}
%\newproof{pot}{Proof of Theorem \ref{thm}}

% put this after all packages
% \makeatletter
% \def\fps@figure{htbp}
% \def\fps@table{htbp}
% \makeatother

\begin{document}

\let\WriteBookmarks\relax

\def\floatpagepagefraction{1}
\def\textpagefraction{.001}

% Short title
\shorttitle{XXX}

% Short author
\shortauthors{K Horikoshi et al.}

% Main title of the paper
\title [mode = title]{XXX}                      
% Title footnote mark
% eg: \tnotemark[1]
% \tnotemark[1,2]

% Title footnote 1.
% eg: \tnotetext[1]{Title footnote text}
% \tnotetext[<tnote number>]{<tnote text>} 
% \tnotetext[1]{This document is the results of the research
%    project funded by the National Science Foundation.}

% \tnotetext[2]{The second title footnote which is a longer text matter
%    to fill through the whole text width and overflow into
%    another line in the footnotes area of the first page.}

\author[1]{Kazuki Horikoshi}[type=editor,
                        auid=000,bioid=1,
                        % role=Researcher,
                        orcid=0009-0003-0521-3099]

% Footnote of the first author
%\fnmark[1]
% \fntext[fn1]{This is the first author footnote. but is common to third author as well.}

% Email id of the first author
\ead{hkazuki@u.nus.edu}
% URL of the first author
%\ead[url]{www.cvr.cc, cvr@sayahna.org}

%  Credit authorship
\credit{Writing – review & editing, Writing – original draft, Visualization, Methodology, Investigation, Analysis, Data curation, Conceptualization.}

% Address/affiliation
\affiliation[1]{organization={Department of the Built Environment, National University of Singapore},
    %addressline={}, 
    city={Singapore},
    % citysep={}, % Uncomment if no comma needed between city and postcode
    postcode={117566}, 
    % state={},
    country={Singapore}}

% Third author
\author[1]{Adrian Chong}[%
   % role=Co-ordinator,
   ]
% Corresponding author indication
\cormark[1]

%\fnmark[3]
\ead{adrian.chong@nus.edu.sg}
%\ead[URL]{www.xxx.org}

\credit{Writing – review & editing, Supervision, Methodology, Conceptualization}

% Corresponding author text
\cortext[cor1]{Corresponding author}
%\cortext[cor2]{Principal corresponding author}

% For a title note without a number/mark
% \nonumnote{This note has no numbers. In this work we demonstrate $a_b$ the formation Y\_1 of a new type of polariton on the interface between a cuprous oxide slab and a polystyrene micro-sphere placed on the slab.
%   }

\acresetall

% \begin{abstract}

Traditional HVAC systems, designed with a "one-size-fits-all" approach, often fail to provide thermal comfort in modern office environments, which are characterized by increasingly flexible and dynamic usage.
Occupant-centric control (OCC) has emerged as a promising solution to enhance both occupant satisfaction and energy efficiency by personalizing thermal conditions.
However, its widespread real-world adoption is hindered by two critical challenges: (1) the inability of existing frameworks to effectively manage the highly dynamic and unpredictable occupancy patterns typical of post-pandemic workplaces, and (2) the prohibitive survey effort required to build and maintain accurate personal comfort models (PCMs) for every individual.

This research aims to address these two gaps by developing and validating a practical OCC framework designed for real-world, multi-occupant office environments.
The methodology consists of a two-part field experiment conducted in an operational office building.
The first experiment will systematically investigate the relationship between various occupancy patterns—such as occupancy variability and user turnover—and the performance of dynamic OCC.
Its objective is to identify the specific conditions under which real-time, personalized control offers significant advantages over conventional static control strategies.

Building on these findings, the second experiment will propose and evaluate a novel, low-effort framework for the operational development of comfort models.
This framework employs an active transfer learning approach to intelligently minimize the total survey burden by strategically deciding when, where, and from whom to collect comfort feedback.
The system prioritizes queries that are most informative for improving group comfort and unlocking energy-saving potential, thereby balancing data efficiency with operational effectiveness.

The anticipated outcome of this research is a set of data-driven, actionable guidelines for building managers on selecting and deploying appropriate OCC strategies based on workspace usage patterns.
Furthermore, the proposed low-effort framework will provide a scalable and economically viable solution for implementing advanced, personalized control in commercial buildings.
By bridging the gap between academic theory and practical application, this work will contribute significantly to creating comfort and energy-efficient indoor environments tailored to the demands of modern work styles.

\end{abstract}

\begin{abstract}
Occupant-centric HVAC control often assumes fixed room-level comfort, however post-COVID offices have become more sparsely and dynamically occupied, creating a mismatch between thermal comfort modeling and room-level control.
This increases the need for higher-resolution group comfort models supported by real-time occupancy data.
This study investigates how the benefit of higher-resolution group comfort models, including use of occupancy tracking, changes with group size in HVAC control.
To examine this, an agent-based simulation was conducted using one month of anonymized location logs and thermal comfort surveys from office rooms of varying sizes and uses, where each agent represented an occupant proxy carrying an observed office movement record.
Three group comfort models were compared: static, based on the room’s regular members; daily, based on that day’s attendees; and real-time, based on the occupants present at each time step.
Personal comfort models (PCM) were developed for 39 participants.
To generate multiple cases, occupancy patterns were created by repeatedly sampling sub-groups from each room’s member pool, and PCMs were randomly assigned to those occupant proxies.
Their “no change” thermal preference probability curves were then merged to form each group model.
Setpoints were selected to maximize the probability of “no change,” and improvement in predicted comfort probability was evaluated against a fixed 24°C baseline.
The results show that, relative to the 24°C baseline, the static model’s improvement falls to nearly zero as room size approaches six occupants, whereas the daily and real-time models maintain improvement in larger rooms.
Further analysis shows that the real-time model’s comfort improvement is greatest under sparse occupancy and becomes limited as mean occupancy approaches four to five occupants.
These findings indicate that higher-resolution comfort aggregation is most beneficial when occupancy remains sparse, and identify when high-resolution occupancy tracking is justified for practical HVAC control.
\end{abstract}

% Use if graphical abstract is present
% \begin{graphicalabstract}
% \includegraphics{figs/grabs.pdf}
% \end{graphicalabstract}

% Research highlights
\begin{highlights}
\item XXX
\item XXX
\end{highlights}

% Keywords Each keyword is seperated by \sep
\begin{keywords}
XXX \sep XXX \sep
\end{keywords}

\maketitle

\section*{Nomenclature}
\renewcommand{\baselinestretch}{0.75}\normalsize
\renewcommand{\aclabelfont}[1]{\textsc{\acsfont{#1}}}
\begin{acronym}[longest]

    \acro{tdb}[$t_{db}$]{dry-bulb air temperature\acroextra{, $^{\circ}$C}}
    \acro{twb}[$t_{wb}$]{wet-bulb air temperature\acroextra{, $^{\circ}$C}}
    \acro{ti}[$t_{i}$]{indoor air temperature\acroextra{, $^{\circ}$C}}
    \acro{tout}[$t_{out}$]{outdoor air temperature\acroextra{, $^{\circ}$C}}
    \acro{tos}[$t_{out,switch}$]{switch-over outdoor temperature\acroextra{, $^{\circ}$C}}
    \acro{to}[$t_{o}$]{operative air temperature\acroextra{, $^{\circ}$C}}
    \acro{tcl}[$t_{cl}$]{clothing temperature\acroextra{, $^{\circ}$C}}
    \acro{tg}[$t_{g}$]{globe temperature\acroextra{, $^{\circ}$C}}
    \acro{rh}[$RH$]{relative humidity\acroextra{, \%}}
    \acro{v}[$V$]{measured air speed\acroextra{, m/s}}
    \acro{vr}[$V_{r}$]{relative air speed\acroextra{, m/s}}
    \acro{vo}[$V_{o}$]{outdoor wind speed\acroextra{, m/s}}
    \acro{tr}[$\overline{t_{r}}$]{mean radiant temperature\acroextra{, $^{\circ}$C}}
    \acro{clo}[$I_{cl}$]{intrinsic clothing insulation from the skin to the outer surface under reference conditions\acroextra{, clo}}
    \acro{clor}[$I_{cl,r}$]{actual intrinsic clothing insulation under given environmental and activities\acroextra{, clo}}
    \acro{met}[$M$]{rate of metabolic heat production\acroextra{, W/m\textsuperscript{2}}}
    \acro{work}[$W$]{rate of mechanical work accomplished\acroextra{, W/m\textsuperscript{2}}}
    
    \acro{pmv}[PMV]{Predicted Mean Vote}
    \acro{ppd}[PPD]{Predicted Percentage of Dissatisfied\acroextra{, \%}}
    \acro{set}[SET]{Standard Effective Temperature\acroextra{, $^{\circ}$C}}
    \acro{tsv}[TSV]{Thermal Sensation Vote}
    \acro{tpv}[TPV]{Thermal Preference Vote}
    \acro{tcv}[TCV]{Thermal Comfort Survey}

    \acro{occ}[OCC]{Occupant-centric control}
    \acro{abw}[ABW]{Activity-based working}
    \acro{bms}[BMS]{building management system}
    \acro{pcm}[PCM]{personal comfort model}
    \acro{gcm}[GCM]{group comfort model}
    \acro{utr}[UTR]{user turnover rate}
    \acro{hvac}[HVAC]{heating, ventilation, and air-conditioning}

    \acro{sgcm}[SGCM]{static group comfort model}
    \acro{dgcm}[DGCM]{daily group comfort model}
    \acro{rtgcm}[RT-GCM]{real-time group comfort model}

    \acro{dml}[DML]{double machine learning}
    \acro{logit}[Logit]{logistic regression}
    \acro{nn}[NN]{neural network}
    \acro{rf}[RF]{random forest}
    % \acro{xgb}[XGB]{extreme gradient boosting}
    \acro{pr2}[partial $R^2$]{partial coefficient of determination}
    \acro{ic}[IC]{Increment accuracy}


\end{acronym}
\renewcommand{\baselinestretch}{1}\normalsize

\section{Introduction}

\SetPicSubDir{ch1-Intro}
\label{sec:Intro}

Thermal comfort is a critical factor influencing not only occupant satisfaction but also productivity and health in indoor environments \cite{bueno_evaluating_2021,kaushik_effect_2022}. 
Studies have increasingly demonstrated the link between indoor environmental quality and occupant outcomes. 
For instance, \mycite{bueno_evaluating_2021} \cite{bueno_evaluating_2021} confirmed this link, reporting that a substantial body of research demonstrates a significant relationship between indoor environmental conditions and occupant productivity.

Traditional \ac{hvac} systems often rely on static control strategies, typically aiming to maintain a uniform temperature based on pre-defined setpoints or the \ac{pmv} model \cite{fanger_thermal_1970}.  
% The PMV model calculates a thermal sensation score using six input variables: two personal factors (clothing insulation and metabolic rate) and four environmental factors (air temperature, mean radiant temperature, air speed, and relative humidity).
While \ac{pmv} attempts to predict the average occupant's thermal sensation, it fundamentally neglects individual differences in thermal preferences \cite{wang_individual_2018}.  
Studies consistently demonstrate that even when \ac{pmv} predicts comfort, a significant percentage of occupants experience discomfort due to these variations \cite{cheung_analysis_2019}.  
This one-size-fits-all approach is inherently limited in achieving widespread occupant satisfaction.
This limitation reflects a difference in prediction target: the \ac{pmv} model represents the mean response of a population, whereas a \ac{pcm} aims to predict the response of a specific individual.
Reviews and field-data analyses have documented substantial between-person variation and reduced predictive accuracy when population models are applied to individuals \cite{wang_individual_2018,cheung_analysis_2019,kim_personal_2018-1}.
These \acp{pcm} address this limitation by learning an individual's response from personal feedback and contextual data, although their data requirements, updating procedures, and integration with building controls remain application-dependent \cite{kim_personal_2018-1}.

\subsection{Occupant-centric control and shared-space comfort aggregation}

In recent years, \ac{occ} has emerged as a promising \ac{hvac} strategy that uses occupant information to support more individualized building operation \cite{soleimanijavid_challenges_2024}.
IEA EBC Annex 79 defines \ac{occ} as an ``occupant-in-the-loop'' approach that seeks to provide optimal and personalized conditions rather than imposing fixed settings \cite{wagner_international_2024}.
In this context, occupant information can range from presence and count to activity and identity, and its usefulness depends on both spatial and temporal resolution \cite{nagy_ten_2023}. Higher information resolution can enable more individualized control, but it also increases sensing, data-management, integration, and privacy requirements.
Typical \ac{occ} frameworks consist of three connected layers: the sensing layer collects information such as environmental conditions, comfort feedback, and occupancy; the learning layer develops preference or occupancy models from these data; and the control layer translates the model outputs into \ac{hvac} actions \cite{soleimanijavid_challenges_2024,xie_review_2020}.
The performance of this chain depends not only on model accuracy but also on sensing quality, communication with the \ac{bms}, controller design, and local operating conditions \cite{zhang_impact_2023}.
Reported barriers to broader deployment include data availability and management, interoperability with existing systems, operator capacity, scalability, privacy, and validation across buildings and climates \cite{obrien_key_2020,nagy_ten_2023,soleimanijavid_challenges_2024}.
Among these functions, the connection between personal comfort information and the control decision is especially important because it determines how individual preferences affect the final setpoint.

In shared office spaces, however, \ac{occ} cannot rely only on isolated \acp{pcm}, because a single room-level \ac{hvac} setpoint must represent multiple occupants at the same time.
Earlier work on shared offices noted that responding directly to one person's feedback implicitly assumes either single occupancy or identical preferences among occupants \cite{schumann_predicting_2010}.
More recent collective-comfort studies likewise frame shared-zone control as a problem of combining heterogeneous preferences under a common environmental decision \cite{topak_collective_2023}.
This makes the construction of a \ac{gcm}, which aggregates individual comfort information into a room-level representation, a central design issue for shared-space \ac{occ}.
The key question is therefore not only how accurately individual \acp{pcm} predict comfort, but also which occupants should be represented when a room-level control decision is made.
Despite growing recognition in the industry and active framework development in simulations, real-world \ac{occ} applications in buildings remain limited~\cite{soleimanijavid_challenges_2024}.

\subsection{Dynamic occupancy and real-time tracking in activity-based working environments}

Meanwhile, the rise of remote and hybrid work has increased the use of \ac{abw} and other flexible office arrangements~\cite{marzban_review_2023}.
In such workplaces, occupants arrive, leave, and move between different spaces throughout the day, causing fluctuations in both occupancy density and occupant composition within a zone \cite{motuziene_office_2022,pan_environmental_2025,huang_space-matching_2025,chang_statistical_2013,levene_problem_2020}.
Dynamic occupancy should therefore be understood not only as a change in the number of people in a room, but also as a change in who is present.
This distinction corresponds to different grades of occupant information: occupancy count describes how many people are present, whereas identity-level information is needed to associate the present group with person-specific comfort profiles \cite{nagy_ten_2023}.
Count data alone therefore cannot determine whether the group comfort representation should change when one occupant replaces another.
Even when the occupancy count is similar, the thermal preferences represented by the present group may differ.
This makes static temperature control less reliable, because a condition optimized for one attendance pattern may not represent another.

% To incorporate occupancy information in \ac{occ} frameworks, two main approaches can be identified: prediction-based methods and real-time sensing.
% Many studies have attempted to predict occupancy patterns using historical data and machine learning techniques.
% However, as noted in prior research \cite{goyal_occupancy-based_2013}, such prediction methods often entail high computational costs, may be more expensive than real-time sensing, and may not offer sufficient accuracy or reliability.

Recent IoT and building-management technologies have made room- or occupant-level occupancy data more accessible through sources such as PIR, CO$_2$, and Wi-Fi-based sensing \cite{soleimanijavid_challenges_2024,zafari_survey_2019,brambilla_potential_2021}.

These modalities do not provide equivalent information.
Reviews identify trade-offs among counting capability, spatial resolution, latency, cost, and privacy for PIR, CO$_2$, camera, WLAN, and related approaches \cite{yang_review_2016,zafari_survey_2019,brambilla_potential_2021}. 
For example, CO$_2$-based inference can respond slowly to changes, while Wi-Fi-based inference can be affected by unstable signals and device-to-occupant behaviour \cite{wang_modeling_2017}.
Identity-level tracking introduces an additional privacy and data-governance burden \cite{obrien_key_2020,nagy_ten_2023}.

\autoref{table:OccuTrack} provides examples of office projects by Japanese companies that use IoT data for individual-level real-time occupancy tracking.
Therefore, opportunities are emerging to use such data in \ac{occ} systems that respond to changes in the occupants currently present.
At the same time, the practical value of real-time occupant-composition data depends on whether the comfort benefit justifies the sensing, integration, privacy, and control burden.

\subsection{Shared-space GCM resolution under dynamic occupancy}

For shared-space \ac{occ}, dynamic occupancy raises a second issue: the controller must decide how individual comfort information should be aggregated for a room-level setpoint.
Existing \ac{occ} research has largely focused on developing and validating \acp{pcm} and control algorithms in single-occupant settings, while shared spaces with multiple occupants remain less explored \cite{park_critical_2019,jung_human---loop_2019,xie_review_2020}.
In shared workspaces, multiple occupants with potentially conflicting thermal preferences must be accommodated simultaneously, and simply averaging preferences or applying a single \ac{pcm} may fail to represent the group.

The difficulty is shaped by both group size and preference heterogeneity.
Shared-space studies have shown that collective comfort depends on how conflicting individual preferences are represented and combined, rather than on occupancy count alone \cite{schumann_predicting_2010,topak_collective_2023}.
Simulation studies further suggest that the combined energy-and-comfort benefit of integrating \acp{pcm} can diminish as more occupants share a thermal zone \cite{jung_energy_2020}, while the potential maximum satisfaction rate may also decrease for larger groups \cite{wang_enhancing_2026}.
% \mycite{jung_energy_2020}\cite{jung_energy_2020} showed that the number of occupants in a zone can affect both comfort and energy performance, and that results depend on how individual comfort models are integrated.
\mycite{ono_effects_2022}\cite{ono_effects_2022} further showed that mismatches between comfort-model resolution and control resolution may reduce control effectiveness and occupant satisfaction.
These findings indicate that shared-space \ac{occ} requires attention to both comfort aggregation and room-level control resolution.

\autoref{table:OccuIntegration} compares field studies that integrate multiple \acp{pcm} within the same control zone.
Most existing multi-occupant studies assume static occupant attendance configurations and do not account for occupants arriving and leaving dynamically \cite{lee_implementation_2019,li_indoor_2019,li_personalized_2017,li_indoor_2023,aguilera_thermal_2019,li_experimental_2021}.
Lei et al.~\cite{lei_practical_2022} addressed dynamic occupancy using multiple reinforcement-learning models for frequent occupant combinations, but such an approach can be data- and compute-intensive and may not cover all possible occupant combinations.
Rather than preparing separate control models for many occupant combinations, a more deployable strategy is to vary the temporal resolution at which individual comfort information is aggregated.
This leads to a practical question: when is static, daily, or real-time aggregation of individual comfort information sufficient for room-level control?

\begin{table}[pos=!htbp]
  \centering
  \begin{tabular}{|l|c|c|c|c|c|c|}
\hline
\textbf{Paper} & \textbf{Year} & \textbf{Resolution} & \textbf{Per unit} & \textbf{Algorithm} & \textbf{Occupancy adaption} \\
\hline
\cite{lei_practical_2022} & 2022 & Group & 4 people& RL & Dynamic  \\
\cite{lee_implementation_2019} & 2019 & Group & 2 people & MPC & Static \\
\cite{li_indoor_2019} & 2019 & Room & 4 people & MPC & Static \\
\cite{li_personalized_2017} & 2017 & Room & 7people & RB & Static \\
\cite{li_indoor_2023} & 2023 & Room & 4 people & RL & Static \\
\cite{aguilera_thermal_2019} & 2019 & Room & 16 people & RB & Static \\
\cite{li_experimental_2021} & 2021 & Room & 6 people & RB & Static \\
\hline
\end{tabular}
  \caption{Field OCC studies integrating multiple personal comfort models}
  \label{table:OccuIntegration}
\end{table}

\subsection{Focus of the research}
\label{subsec:Focus}
Based on this background, this study evaluates when higher-resolution aggregation of individual comfort information improves predicted comfort in shared offices, and what operational burden is introduced by real-time occupant-composition updates.
The analysis compares three levels of group-comfort aggregation, from fixed room-member aggregation to daily and real-time occupant-composition updates, with the operational definitions provided in the Method section.
Building on prior work on group size and control granularity \cite{ono_effects_2022,jung_energy_2020}, the analysis introduces temporal occupancy dynamics, especially short-horizon occupant turnover, as a key dimension of shared-space \ac{occ}.
% Using field data from real office spaces, the study relates model performance to subgroup size, mean occupancy, occupancy variability, and \ac{utr}.

Specifically, the analysis compares \ac{sgcm}, \ac{dgcm}, and \ac{rtgcm} policies in terms of:
(1) mean comfort probability and improvement relative to a fixed \qty{24}{\celsius} baseline;
(2) the setpoint adjustment magnitude required by the \ac{rtgcm}; and
(3) the relationship between occupancy dynamics, especially \ac{utr}, and the frequency of actionable control updates.

\FloatBarrier

\section{Method}
\SetPicSubDir{ch3-method}
\label{sec:Method}
This study explores the occupant comfort performance of \ac{occ} employing three types of \ac{gcm} based on different occupancy resolution and the effect of occupancy dynamics on the control.
To examine these purposes, agentic-simulation framework is used based on field data collected in a real-world office building.
The overall procedural flow of the study is illustrated in \autoref{fig:ex_flow}.
\begin{figure}[pos=htbp!]
    \centering
    \includegraphics[width=0.97\linewidth]{\Pic{png}{procedure}}
    \caption{Overall procedure flow of the study}
    \label{fig:ex_flow}
\end{figure}

\subsection{Data Collection}
\label{subsec:DataCollect}
For the simulation, three primary types of field data were collected in the target building: real-time occupancy logs, subjective \ac{tcv} surveys and environmental measurement logs.

\subsubsection{Target Building, Participant Recruitment and Ethical Considerations}

\label{sec:TargetBuil}
The data collection was conducted at The GEAR, an R\&D and administrative office building located in Singapore and operated by Kajima Corporation, from 29 September to 24 October 2025 (4 weeks).
The 3rd to 6th floors, which accommodate a mix of shared-office and static seating work styles depending on departmental roles, were the target areas for data collection.
Data collection was limited to the following rooms: 3F: Admin Office A, Admin Office B; 5F: R\&D Office, Admin Office C, Shared Office 1, Shared Office 2; 6F: Design Office.
% Data collection was limited to the following rooms: 3F: Office 3A, Office 3B; 5F: Katris Office, KD Office, Shared Office 1, Shared Office 2; 6F: ILYA Office.
The target zones are served by a central air-conditioning system consisting of air-handling units and dedicated outdoor air-handling units, using variable air volume units capable of individual zone temperature and ventilation control.
\qty{45} office workers who regularly work in the designated target rooms were invited to participate.
Considering the completeness of the collected survey answers and location tracking records, data from \qty{39} participants were used in the analysis.
This study has been approved by the National University of Singapore Institutional Review Board (NUS-IRB Reference Code: NUS-IRB-2025-498).
Informed consent was obtained from all participants prior to the study.
All location and survey data were anonymized by the building administrator before being accessed by the research team to protect privacy. 


\subsubsection{Thermal Comfort Survey}
Subjective \ac{tcv} survey answers were collected via a web survey during work hours. 
Items and scales follow ANSI/ASHRAE Standard 55:2023, Appendix~L \cite{ashrae_ansiashrae_2022}.
The survey contents are summarized in \autoref{table:comfsurvey}.
Responses within 20 minutes of room entry were excluded to avoid thermal transient effects.
Participants were prompted up to six times per workday: three times in the morning and three times in the afternoon.
Over the one-month observation period, \qty{2608} survey responses were collected in total.

\subsubsection{Indoor Environmental Measurements}
Indoor environmental data was monitored in parallel with \ac{tcv} survey.
The primary measurements were \ac{ti} and \ac{rh}, obtained from compact sensors (HIOKI E.E. Corporation, Nagano, Japan) installed in the target rooms.
For rooms with two to three compact sensors, their mean value was used as a representative room-level value to reduce sensor-specific measurement bias.
To ensure sufficient variability of survey answers, \ac{ti} setpoints were cycled between \qtyrange{22}{27}{\celsius} in \qty{1}{\celsius} steps, with a randomized order across days.  
On selected days, morning and afternoon setpoints differed to expand diversity.
In addition, measured air speed (\acs{v}) and mean radiant temperature (\acs{tr}) were recorded periodically as diagnostic checks rather than as main focus of the analysis.
% TO DO measurement equipment ONDOTORI

\subsubsection{Individual-level location tracking log}
Individual-level location data were obtained from anonymized, timestamped records provided through the \ac{bms}.
The raw zone-level records were reconstructed into stay logs, where each entry represents a continuous period of presence by one participant in one room.
Each stay-log entry contains an anonymized participant ID, room ID, start time, and end time.
These stay logs provide the basis for constructing location-log agents in the simulation.
The detailed procedure for this reconstruction is described in Appendix~\ref{app:occupancy_reconstruction}.
% TO DO: refine output to show as appendix, might be connected to how to show chapters in appendidices

\subsection{Occupancy-log-based agentic comfort simulation}
\subsubsection{Simulation Design}
Using the collected \ac{tcv} responses, environmental measurements, and reconstructed stay logs, two agentic components were prepared for each participant: a \ac{pcm} and a location-log agent.
The \ac{pcm} represents the participant's thermal preference response to indoor temperature, whereas the location-log agent preserves the participant's observed room-level presence pattern over the study period.
In the simulation, these two components were intentionally separated so that comfort-response characteristics could be randomly reassigned to observed occupancy patterns.
A Monte Carlo-style resampling simulation was then conducted by combining sampled location-log agents with randomly assigned \acp{pcm}.
% At each control timestep, the location-log agents determined the simulated room occupant composition.
% Based on the simulated occupant compositions, three \ac{gcm}-based \ac{occ} policies were evaluated from comfort and control perspectives.
% TO DO: mention general room regular user agents are used

\subsubsection{Personal Comfort Model (PCM)}
Individual \acp{pcm} are developed for each participant to predict thermal preference based on room-level \ac{ti}.
The \ac{pcm}s are developed using ordered multinominal logistic classification, used in similar agentic-based thermal comfort simulations~\cite{jung_comparative_2019, jung_energy_2020}.
The model predicts the probability distribution over three thermal preference classes: $S_{th} \in \{+1 (\text{Warmer}), 0 (\text{No change}), -1 (\text{Cooler})\}$, using \ac{ti} as the primary feature.
\begin{equation}
P(S_{th} = s \mid t_i)\,.
\label{eq:pcm_probability}
\end{equation}

% TO DO: need to consider if ref is needed for multinominal logistic classification

\subsubsection{Monte Carlo simulation and GCM policies}
As introduced in the research focus, the simulation evaluates three temporal resolutions of \ac{gcm} aggregation: \ac{sgcm}, based on the sampled regular room-member group over the study period; \ac{dgcm}, based on the sampled attendees on each day; and \ac{rtgcm}, based on the sampled occupants present at each control timestep.

\begin{figure}[pos=!bt]
    \centering
    \includegraphics[width=0.97\linewidth]{\Pic{png}{OverallFlow}}
    \caption{Overall flow of simulation. Step1 (a)From room user pool, a subgroup of location-log agents is sampled and PCMs are randomly assigned to the sampled location-log agents. Step2(c) Based on the assigned PCMs and location logs,(d) Possibility curve of "No change" in thermal preference id merged based on three types of GCM at each control timestep (d) Setpoints are selected at the point where the aggregated comfort probability is maximized (e) The setpoints are evaluated by returning to the assigned PCMs and extracting the predicted ``No change'' probability from each PCMs}
    \label{fig:ComChara}
\end{figure}

Using the prepared location-log agents and \acp{pcm}, Monte Carlo trials were generated to represent different combinations of observed occupancy compositions and assigned thermal-preference profiles for each target room.
To preserve room-specific usage patterns while varying the simulated occupant mix, each Monte Carlo trial sampled a subgroup of location-log agents from the regular member pool of room $r$ and randomly assigned \acp{pcm} to the selected location-log agents.
The subgroup size $K$ denotes the number of location-log agents sampled from the regular member pool of room $r$ in a simulation condition.
Within a trial, the sampled subgroup size $K$ is fixed for the full simulation period.
However, only a subset of those sampled agents is present at any control timestep, so the realized occupancy count varies according to the observed stay logs and is denoted by $N_{r,t}=|S_{r,t}|$.
The selected location-log agents retained their observed stay logs in their original rooms, while \acp{pcm} were randomly assigned to those selected location-log agents in each trial.
This procedure preserved observed room usage patterns while varying the association between occupancy patterns and comfort-response characteristics.
At each 30-minute control timestep $t$, the subset of sampled location-log agents who are present in room $r$ defined the active set $S_{r,t}$, and each policy selected the corresponding temperature setpoint from its \ac{gcm}.

For each Monte Carlo trial, the three \acp{gcm} described above were formed by averaging the assigned individual ``No change'' probability curves at different temporal resolutions.
For each policy, the selected setpoint was the temperature that maximized the corresponding aggregated ``No change'' probability.
Thus, the three policies share the same setpoint-selection rule and differ only in the update timescale of the represented group: the full study period, each day, or each control timestep.
Let $p_i(T)=P_i(S_{th}=0\mid T)$ denote the ``No change'' probability predicted by the \ac{pcm} assigned to selected location-log agent $i$, where $i$ indexes an individual occupant proxy in the simulation, at \acl{ti} $T$.
For a group of \ac{gcm} $G$, the merged group comfort probability is defined as
\begin{equation}
\bar{p}_{G}(T)
=
\frac{1}{|G|}
\sum_{i\in G} p_i(T).
\label{eq:gcm_probability_curve}
\end{equation}
The group-level setpoint is then selected as
\begin{equation}
T_G^{*}
=
\operatorname*{arg\,max}_{T\in\mathcal{T}_{\mathrm{set}}}
\bar{p}_{G}(T),
\label{eq:gcm_optimal_setpoint}
\end{equation}
where $\mathcal{T}_{\mathrm{set}}$ is the temperature grid of candidate setpoints.
The policy-specific groups are
\begin{equation}
G_{r}^{\mathrm{Static}} = R_r,
\qquad
G_{r,d}^{\mathrm{Daily}} = A_{r,d},
\qquad
G_{r,t}^{\mathrm{RT}} = S_{r,t},
\label{eq:policy_groups}
\end{equation}
where $R_r$ is the sampled room member group for room $r$, $A_{r,d}$ is the set of sampled location-log agents who appear in room $r$ on day $d$, and $S_{r,t}$ is the set of sampled location-log agents present at timestamp $t$.
The resulting policy setpoints are
\begin{equation}
T_{\mathrm{set}}^{\mathrm{Static}}(t)=T_{G_{r}^{\mathrm{Static}}}^{*},
\qquad
T_{\mathrm{set}}^{\mathrm{Daily}}(t)=T_{G_{r,d(t)}^{\mathrm{Daily}}}^{*},
\qquad
T_{\mathrm{set}}^{\mathrm{RT}}(t)=T_{G_{r,t}^{\mathrm{RT}}}^{*},
\label{eq:policy_setpoints}
\end{equation}
where $d(t)$ is the calendar day containing timestamp $t$.

In addition to these three group-model-based policies, a fixed \qty{24}{\celsius} setpoint is used as an additional baseline.
For each room and subgroup size, \qty{1000} Monte Carlo trials were formed by combining sampled location-log agents and assigned \acp{pcm}, and one-month \ac{occ} simulations based on the three \acp{gcm} were conducted.


\subsection{Evaluation Metrics}
\label{sec:EvalMet}

\subsubsection{Occupancy-related evaluation parameters}
\label{sec:OccuMetrics}
\autoref{table:OccuMetrics} in~\autoref{app:OccupancyStudy} summarizes occupancy-related literature and metrics after COVID-19.
While many studies emphasize occupancy count and related parameters, this study also considers occupant-composition turnover.
Let $\mathcal{T}_{r,d}$ denote the ordered set of analysis timestamps for room $r$ on day $d$, and let $S_{r,t}$ denote the set of sampled location-log agents present in room $r$ at timestamp $t$.
The realized occupancy count is $N_{r,t}=|S_{r,t}|$.

The subgroup size $K$ is the number of sampled location-log agents in a simulation condition, whereas $N_{r,t}$ is the number of those location-log agents actually present at timestamp $t$.

\paragraph{Realized occupancy count}
\begin{equation}
N_{r,t} \;=\; |S_{r,t}|.
\label{eq:Nt}
\end{equation}

\paragraph{Daily mean occupancy}
The daily mean occupancy is
\begin{equation}
\bar{N}_{r,d}
=
\frac{1}{|\mathcal{T}_{r,d}|}
\sum_{t\in\mathcal{T}_{r,d}} N_{r,t}.
\label{eq:mean_occ_daily}
\end{equation}

% \paragraph{Occupancy variability}
% To capture within-day fluctuations in room load, the daily standard deviation of realized occupancy count is used:
% \begin{equation}
% \sigma_{N,r,d}
% =
% \sqrt{
% \frac{1}{|\mathcal{T}_{r,d}|-1}
% \sum_{t\in\mathcal{T}_{r,d}}
% \left(N_{r,t}-\bar{N}_{r,d}\right)^2
% }.
% \label{eq:occ_var}
% \end{equation}

\paragraph{User turnover rate}
The short-horizon replacement of attendees is quantified using a rolling Jaccard turnover, denoted as \ac{utr}.
For timestamp-level analysis, \ac{utr} is computed by comparing the current occupant set with the set observed 30 minutes earlier:
\begin{equation}
\mathrm{UTR}_{r,t}^{30}
=
1
-
\frac{
|S_{r,t}\cap S_{r,t-30}|
}{
|S_{r,t}\cup S_{r,t-30}|
}.
\label{eq:rolling_turnover}
\end{equation}
Here, $S_{r,t-30}$ denotes the occupant set in the same room 30 minutes before timestamp $t$.
When both the current and lagged sets are empty, the timestamp-level turnover is treated as undefined.
Higher $\mathrm{UTR}_{r,t}^{30}$ indicates stronger short-horizon replacement of attendees.

% For daily summaries, timestamp-level \ac{utr} values are averaged within each room-day:
% \begin{equation}
% \mathrm{UTR}_{r,d}
% =
% \frac{1}{|\mathcal{V}_{r,d}|}
% \sum_{t\in\mathcal{V}_{r,d}}
% \mathrm{UTR}_{r,t}^{30},
% \label{eq:utr_daily}
% \end{equation}
% where $\mathcal{V}_{r,d}$ is the set of timestamps on day $d$ for which the 30-minute turnover is defined.
% The daily \ac{utr} captures how much the present group composition changes over short horizons, even when the occupancy count is similar.


\subsubsection{Comfort-related outcomes}
\label{subsec:comfMetrics}

\paragraph{Policy setpoint}
Let $T_{\mathrm{set}}^{\mathrm{policy}}(t)$ denote the setpoint at time $t$ generated by $\mathrm{policy}\in\{\mathrm{Static},\mathrm{Daily},\mathrm{RT}\}$.
Here, $\mathrm{policy}\in\{\mathrm{Static},\mathrm{Daily},\mathrm{RT}\}$ corresponds to the \ac{sgcm}, \ac{dgcm}, and \ac{rtgcm} policies, respectively.

\paragraph{Occupant-level comfort probability}
For each present member $i\in S_{r,t}$, the setpoint generated by a policy is evaluated by the predicted ``No change'' probability $p_i(T_{\mathrm{set}}^{\mathrm{policy}}(t))$ from the assigned \ac{pcm}.

\paragraph{Room-day comfort probability}
The room-day comfort probability is then computed as the person-time average:
\begin{equation}
\mathrm{CP}_{r,d}^{\mathrm{policy}}
=
\frac{
\sum_{t\in\mathcal{T}_{r,d}}
\sum_{i\in S_{r,t}}
p_i\!\left(T_{\mathrm{set}}^{\mathrm{policy}}(t)\right)
}{
\sum_{t\in\mathcal{T}_{r,d}}
|S_{r,t}|
}.
\label{eq:comfort_probability_daily}
\end{equation}
Here, $\mathcal{T}_{r,d}$ is the ordered set of analysis timestamps for room $r$ on day $d$.
Timestamps with no present location-log agents do not contribute person-time comfort scores.

\paragraph{Mean comfort probability}
For comparison across the simulation period, daily room-level comfort probabilities were averaged over the analysis days and Monte Carlo trials for each room, subgroup size, and policy:
\begin{equation}
\overline{\mathrm{CP}}_{r,K}^{\mathrm{policy}}
=
\frac{1}{|\mathcal{M}_{r,K}|}
\sum_{m\in\mathcal{M}_{r,K}}
\frac{1}{|\mathcal{D}_{r,m}|}
\sum_{d\in\mathcal{D}_{r,m}}
\mathrm{CP}_{r,d,m}^{\mathrm{policy}} .
\end{equation}
where $m$ is the Monte Carlo trial index, $\mathcal{M}_{r,K}$ is the set of Monte Carlo trials for room $r$ and subgroup size $K$, $\mathcal{D}_{r,m}$ is the set of valid analysis days for room $r$ in trial $m$, and $\mathrm{policy}\in\{\mathrm{Static},\mathrm{Daily},\mathrm{RT},24\,^\circ\mathrm{C}\}$ denotes the evaluated control policy.

\paragraph{Comfort probability improvement}
The improvement relative to the fixed \qty{24}{\celsius} baseline is computed as
\begin{equation}
\Delta \overline{\mathrm{CP}}_{r,K}^{\mathrm{policy}}
=
\overline{\mathrm{CP}}_{r,K}^{\mathrm{policy}}
-
\overline{\mathrm{CP}}_{r,K}^{24\,^\circ\mathrm{C}} .
\label{eq:comfort_probability_improvement}
\end{equation}

\subsubsection{Control-related outcomes}

\paragraph{Setpoint adjustment magnitude}
To quantify setpoint adjustment requirements under the real-time policy, the setpoint adjustment from the previous control timestep is defined as
\begin{equation}
\delta T_{r,t}^{\mathrm{RT}}
=
T_{\mathrm{set}}^{\mathrm{RT}}(t)
-
T_{\mathrm{set}}^{\mathrm{RT}}(t-\Delta t),
\label{eq:rt_setpoint_step}
\end{equation}
where $\Delta t=\qty{30}{\minute}$ is the control timestep, and $\delta T_{r,t}^{\mathrm{RT}}$ is the setpoint change required by the real-time policy from the previous control timestep.
The adjustment magnitude is evaluated as $|\delta T_{r,t}^{\mathrm{RT}}|$.
Because this analysis focuses on actual setpoint adjustments, the daily mean adjustment magnitude is computed only over timesteps with a non-negligible setpoint change.
\begin{equation}
\mathcal{C}_{r,d}^{+}
=
\left\{
t\in\mathcal{C}_{r,d}
\;:\;
\left|\delta T_{r,t}^{\mathrm{RT}}\right| \geq \varepsilon
\right\},
\qquad
\varepsilon=\qty{0.05}{\celsius}.
\label{eq:changed_step_set}
\end{equation}
\begin{equation}
\overline{|\delta T|}_{r,d,+}^{\mathrm{RT}}
=
\frac{1}{|\mathcal{C}_{r,d}^{+}|}
\sum_{t\in\mathcal{C}_{r,d}^{+}}
\left|\delta T_{r,t}^{\mathrm{RT}}\right|,
\label{eq:rt_mean_abs_step_when_changed_daily}
\end{equation}
where $\mathcal{C}_{r,d}$ is the set of timestamps on day $d$ in room $r$ for which both $T_{\mathrm{set}}^{\mathrm{RT}}(t)$ and $T_{\mathrm{set}}^{\mathrm{RT}}(t-\Delta t)$ are defined, $\mathcal{C}_{r,d}^{+}$ is the subset of valid control timesteps with a non-negligible setpoint change, and $\overline{|\delta T|}_{r,d,+}^{\mathrm{RT}}$ is the mean absolute adjustment magnitude when the setpoint changes.

\paragraph{Action-needed window ratio}
For the Results analysis of control action frequency, the main actionable metric is the 30-minute action-needed window ratio.
For the 30-minute decision-window analysis, the real-time setpoint is first rounded to the nearest \qty{0.5}{\celsius}, denoted by $T_{\mathrm{set},0.5}^{\mathrm{RT}}(t)$.
A 30-minute window is counted as requiring action when at least one rounded real-time setpoint update within the window is \qty{0.5}{\celsius} or larger:
\begin{equation}
\mathrm{AWR}_{r,d}^{30}
=
\frac{1}{|\mathcal{W}_{r,d}^{30}|}
\sum_{w\in\mathcal{W}_{r,d}^{30}}
\mathbb{1}
\left(
\max_{t\in\mathcal{C}_{r,d,w}}
\left|
T_{\mathrm{set},0.5}^{\mathrm{RT}}(t)
-
T_{\mathrm{set},0.5}^{\mathrm{RT}}(t-\Delta t)
\right|
\geq \qty{0.5}{\celsius}
\right),
\label{eq:action_needed_window_ratio}
\end{equation}
where $\mathbb{1}(\cdot)$ is the indicator function, $\mathcal{W}_{r,d}^{30}$ is the set of 30-minute decision windows for room $r$ on day $d$, $\mathcal{C}_{r,d,w}$ is the set of valid within-window consecutive setpoint-update timestamps in window $w$, and $\mathrm{AWR}_{r,d}^{30}$ is the share of 30-minute windows requiring at least one actionable real-time setpoint update.

The defined metrics are used for three comparisons: mean comfort probability and improvement relative to the fixed \qty{24}{\celsius} baseline across subgroup sizes and \ac{gcm} resolutions; real-time setpoint adjustment requirements through mean absolute setpoint adjustment; and the relationship between \ac{utr} and the 30-minute action-needed window ratio.
Representative \ac{utr} transitions by room and weekday are then used to interpret the occupancy-dynamics patterns associated with higher action-needed ratios.

% Removed

% \subsubsection{Personal Comfort Model (PCM)}

% Individual \acp{pcm} are developed for each participant to predict thermal preference based on indoor air temperature.
% Environmental variables are sourced from compact indoor environment sensors (indoor air temperature (\acs{ti}).
% The baseline \ac{pcm} are a multi-class logistic classification model, used in similar agentic-based thermal comfort simulations~\cite{jung_comparative_2019, jung_energy_2020}.

% %study by Daum et al.~\cite{daum_personalized_2011}, widely adopted in prior studies~\cite{jung_comparative_2019,lee_implementation_2019,zhang_impact_2023,meimand_personal_2024,topak_collective_2023,gupta_incentive-based_2018}.

% Specific model inputs and outputs are defined as follows:
% \begin{itemize}
%     \item \textbf{Input vector $\mathbf{x}$:} The primary feature is indoor air temperature (\acs{ti}) obtained from the compact sensors nearest to the occupant's zone.
%     \item \textbf{Output vector $\mathbf{y}$:} The model predicts the probability distribution over three thermal preference classes: $S_{th} \in \{+1 (\text{Warmer}), 0 (\text{No change}), -1 (\text{Cooler})\}$.
% \end{itemize}

% \begin{equation}
% P(S_{th} = s \mid t_i)\,.
% \label{eq:pcm_probability}
% \end{equation}

% \begin{figure}[pos=!bt]
%     \centering
%     \includegraphics[width=0.5\linewidth]{\Pic{png}{ex_comfortmodel_example}}
%     \caption{Individual comfort models from comfort surveys}
%     \label{fig:ComChara}
% \end{figure}

% From each participant’s \ac{pcm}, individual metrics are derived:
% \begin{description}
%   \item[Preferred Temperature $T^p_i$:]
%   The temperature that maximizes the probability of “no change,” i.e.,
%   \begin{equation}
%     T^p_i = \operatorname*{arg\,max}_{t} P(S_{th}=0 \mid t)\,.
%   \label{eq:preferred_temperature}
%   \end{equation}
%   \item[Individual Comfort Range {\([T^c_i,\,T^h_i]\)}:]
%   The set of temperatures at which the probability of ``no change'' dominates the two preference alternatives (``warmer'' and ``cooler''):
%   \begin{equation}
%   [T^c_i,\,T^h_i] \;\equiv\; \left\{\, t \;\middle|\; P(S_{th}=0 \mid t)\;\ge\;\max\!\bigl(P(S_{th}=+1\mid t),\,P(S_{th}=-1\mid t)\bigr) \,\right\}.
%   \label{eq:individual_comfort_range}
%   \end{equation}
%   This range will be used to evaluate how frequently a given setpoint schedule places present occupants within their personal comfort ranges.
% \end{description}


% Previous version retained for checking:
% The corresponding group-level setpoints are defined as
% \[
% T_{\mathrm{static}} \;=\; \frac{1}{|R|}\sum_{i\in R}\theta_i,
% \qquad
% T_{\mathrm{daily}}(d) \;=\; \frac{1}{|A_d|}\sum_{i\in A_d}\theta_i,
% \qquad
% \mu_t \;=\; \frac{1}{|S_t|}\sum_{i\in S_t}\theta_i \;\equiv\; T_{\mathrm{opt,GCM}}(t).
% \]
% Here, $T_{\mathrm{static}}$, $T_{\mathrm{daily}}(d)$, and $\mu_t$ correspond to the \ac{sgcm}, \ac{dgcm}, and \ac{rtgcm}, respectively.

% Previous binary coverage version retained for checking:
% For each occupant $i \in S_t$, define
% \begin{equation}
% C_{i,t}^{\mathrm{policy}} =
% \begin{cases}
% 1, & \text{if } T^c_i \le T_{\mathrm{set}}^{\mathrm{policy}}(t) \le T^h_i,\\[4pt]
% 0, & \text{otherwise.}
% \end{cases}
% \label{eq:coverage_indicator}
% \end{equation}
%
% Room-level coverage over a target period $\mathcal{T}_r$ is
% \begin{equation}
% \mathrm{Coverage}_{r}^{\mathrm{policy}}
% =\frac{\sum_{t\in\mathcal{T}_r}\sum_{i\in S_t} w_{i,t}\, C_{i,t}^{\mathrm{policy}}}{\sum_{t\in\mathcal{T}_r}\sum_{i\in S_t} w_{i,t}},
% \quad w_{i,t}=1\ \text{unless stated otherwise.}
% \label{eq:coverage_room_level}
% \end{equation}


% Previous occupancy metric version retained for checking:
% \paragraph{Occupancy count (target group size effect).}
% \begin{equation}
% N_t \;=\; |S_t|.
% \label{eq:Nt}
% \end{equation}
% Larger groups average out individual preferences more strongly and are expected to reduce $|\Delta T|$.
%
% \paragraph{Occupancy variability (variability of occupancy count).}
% To capture short-horizon fluctuations in room load, a windowed dispersion of counts is used, e.g., for window size $W$,
% \begin{equation}
% \sigma^{(W)}_{N,t} \;=\; \operatorname{SD}\big(\,N_{t-W+1},\,\dots,\,N_t\,\big).
% \label{eq:occ_var}
% \end{equation}
% Greater variability indicates volatile attendance and is expected to increase $|\Delta T|$.
%
% \paragraph{\ac{utr}.}
% The data are quantified in terms of short-horizon replacement of attendees via a rolling Jaccard turnover, denoted as \ac{utr}:
% \begin{equation}
% \mathrm{UTR}^{(W)}_t \;=\; 1 \;-\; \frac{|S_t \cap S_{t-W}|}{|S_t \cup S_{t-W}|}.
% \label{eq:rolling_turnover}
% \end{equation}
% Higher $\mathrm{UTR}^{(W)}_t$ indicates more sustained replacement of attendees and is expected to increase $|\Delta T|$.

% From the stay logs, daily occupancy metrics are defined as:

% \begin{itemize}
%     \item \textbf{Average Occupancy Count ($N_{\text{avg}}$):} 
%     The mean number of occupants per day, where $N_{\text{avg}}$ is calculated as the daily average of all logged stays.

%     \item \textbf{Occupancy Variability ($V_{occ}$):} 
%     The normalized standard deviation of daily occupancy counts, representing how much the number of occupants fluctuates,
%     \[
%       V_{occ} = \frac{\sigma_{\text{log}}}{N_{\text{avg}}}
%     \]
%     where $\sigma_{\text{log}}$ is the standard deviation of the stay log counts.

%     \item \textbf{Occupant Turnover Rate (OTR):} 
%     The ratio of the number of unique occupants appearing in a day to the average occupancy, representing how frequently different individuals rotate in and out of the space,
%     \[
%       \text{OTR} = \frac{N_{\text{unique}}}{N_{\text{avg}}}
%     \]
%     where $N_{\text{unique}}$ is the total number of distinct occupant IDs observed on that day.

%     \item \textbf{Day-to-Day Turnover Rate (DTR):} 
%     The Jaccard distance between occupant sets of consecutive days, indicating how much the attendance composition changes from one day to the next,
%     \[
%       \text{DTR}_d = 1 - \frac{|D_d \cap D_{d-1}|}{|D_d \cup D_{d-1}|}
%     \]
%     where $D_d$ and $D_{d-1}$ denote the sets of occupants present on day $d$ and the previous day. 
%     A higher value indicates a greater turnover in the daily group composition.
% \end{itemize}

% \item For the Daily baseline, larger plan--actual mismatch increases the deviation from the daily setpoint:
% \[
% \mathbb{E}\!\left[\,|\Delta T^{\mathrm{Daily}}_t| \,\big|\, M^{\mathrm{plan}}_t \,\right] 
% \;\;\text{is an increasing function of}\; M^{\mathrm{plan}}_t.
% \]


% Previous control-deviation version retained for checking:
% Among the three \acp{gcm} defined above, the \ac{rtgcm} setpoint $\mu_t$ is treated as the highest-resolution target and is compared with the \ac{sgcm} and \ac{dgcm} baselines.
% The deviations are
% \begin{equation}
% \Delta T^{\mathrm{Static}}_t \;=\; \mu_t - T_{\mathrm{static}},
% \qquad
% \Delta T^{\mathrm{Daily}}_t \;=\; \mu_t - T_{\mathrm{daily}}(d).
% \label{eq:deltaT_defs}
% \end{equation}
% Unless stated otherwise, the magnitudes of $|\Delta T^{\mathrm{Static}}_t|$ and $|\Delta T^{\mathrm{Daily}}_t|$ will be analyze

% \subsubsection{Occupancy-stratified summaries}
% To examine how occupancy conditions relate to coverage, stratified summaries will be reported by discretizing occupancy metrics into ordered bins
% (e.g., quartiles): group size $N_t$, short-horizon count variability $\sigma^{(W)}_{N,t}$,
% and user turnover rate $\mathrm{UTR}^{(W)}_t$ (Section~\ref{sec:OccuMetrics}).
% For a bin $g$, define
% \begin{equation}
% \mathrm{Coverage}_{r,g}^{\mathrm{policy}}
% =\frac{\sum_{t\in\mathcal{T}_r \cap \mathcal{G}_g}\sum_{i\in S_t} C_{i,t}^{\mathrm{policy}}}{\sum_{t\in\mathcal{T}_r \cap \mathcal{G}_g}\sum_{i\in S_t} 1},
% \label{eq:coverage_stratified}
% \end{equation}
% and contrasts
% \begin{align}
% \Delta \mathrm{Cov}_{r,g}^{\mathrm{Dyn-Daily}}
% &= \mathrm{Coverage}_{r,g}^{\mathrm{Dynamic}} - \mathrm{Coverage}_{r,g}^{\mathrm{Daily}},
% \label{eq:coverage_contrast_dyn_daily} \\
% \Delta \mathrm{Cov}_{r,g}^{\mathrm{Dyn-Static}}
% &= \mathrm{Coverage}_{r,g}^{\mathrm{Dynamic}} - \mathrm{Coverage}_{r,g}^{\mathrm{Static}}.
% \label{eq:coverage_contrast_dyn_static}
% \end{align}

% \subsubsection{Statistical Comparison and Equivalence Testing}
% \label{subsec:equivalence}

% It will be assessed whether the Dynamic setpoint is equivalent to each baseline using
% two one-sided tests (TOST) with an equivalence margin of $\pm 0.5^\circ\mathrm{C}$
% (aligned with the general thermostat control step). 
% The null posits non-equivalence; rejection of both one-sided tests indicates equivalence under the specified occupancy conditions.

% \begin{enumerate}
%     \item \textbf{Outcomes}
%     Let $\Delta T^{\mathrm{Static}}_t$ and $\Delta T^{\mathrm{Daily}}_t$ be defined as in \eqref{eq:deltaT_defs}.
%     Unless stated otherwise, the magnitudes of $|\Delta T^{\mathrm{Static}}_t|$ and $|\Delta T^{\mathrm{Daily}}_t|$ will be analyzed.
%     \item \textbf{Hypotheses and margin}
%     For each baseline, we test equivalence to Dynamic with the margin $\pm 0.5^\circ\mathrm{C}$ using TOST.

%     \item \textbf{Procedure}
%     Apply two one-sided tests at the pre-specified significance level; equivalence is concluded if both
%     null hypotheses are rejected.

% \end{enumerate}

% \subsubsection{Energy Impact Estimation under Dynamic GCM}
% \label{subsec:energy_impact}

% The counterfactual energy use will be estimated under each policy by combining actual exogenous conditions, including outdoor air temperature (\acs{tout}), with policy-specific setpoint schedules.

% Let $E_k(t)$ denote an energy or power variable from the \acs{bms} at time $t$ for component $k \in \mathcal{K}$ (e.g., supply fan power, cooling/heating coil energy etc).
% \[
% \mathbf{x}(t) \;=\; \bigl(t_{out}(t),\, T_{\mathrm{set}}(t),\, N_t,\, \ldots \bigr).
% \]
% % GP model
% Each component is modeled via Gaussian Process regression:
% \[
% E_k(t)
% \;=\;
% f_k\!\bigl(\mathbf{x}(t)\bigr) \;+\; \varepsilon_{k,t},
% \qquad k \in \mathcal{K},
% \]

% where $f_k(\cdot)$ is a flexible nonparametric function and $\varepsilon_{k,t}$ captures residual variability.

% Holding actual exogenous series fixed (e.g., $t_{out}(t)$ and $N_t$), the counterfactual energy under each policy is generated by feeding the corresponding setpoint schedule:
% \[
% \widehat{E}_{k}^{\mathrm{policy}}(t)
% \;=\;
% \widehat{f}_k\!\bigl(t_{out}(t),\, T_{\mathrm{set}}^{\mathrm{policy}}(t),\, N_t,\, \ldots \bigr).
% \]
% The value is then aggregated over the study period:
% \[
% \mathrm{TotalEnergy}_{k}^{\mathrm{policy}}
% \;=\;
% \sum_{t \in \mathcal{T}} \widehat{E}_{k}^{\mathrm{policy}}(t),
% \qquad
% \mathrm{SavingRate}_{k}^{\mathrm{Dyn\;vs\;baseline}}
% \;=\;
% 1 - \frac{\mathrm{TotalEnergy}_{k}^{\mathrm{Dynamic}}}{\mathrm{TotalEnergy}_{k}^{\mathrm{baseline}}},
% \]
% where the baseline is either Static or Daily.

% The total amount will be reported by component and overall, along with uncertainty intervals from the GP predictive distribution.


% \paragraph{Modeling strategy and Feature Attribution}

% Because the relationship between occupancy metrics and the absolute deviations cannot be assumed linear, the absolute deviation $|\Delta T|$ will be modeled as a function of occupancy metrics:
% \begin{equation}
% |\Delta T| \;=\; f\!\bigl(N_{\text{avg}},\, V_{\text{occ}},\, \mathrm{UTR}^{(W)}_t\bigr) \;+\; \varepsilon,
% \label{eq:deltaT_gp_model}
% \end{equation}
% where $f(\cdot)$ is approximated using Gaussian Process regression.
% This allows us to capture potentially non-linear relationships between occupancy conditions and the deviation magnitude.

% For interpretability, feature attribution methods such as SHAP (SHapley Additive exPlanations) will be applied to the GP predictions, so that the relative contribution of each occupancy metric ($N_{\text{avg}}, V_{\text{occ}}, \mathrm{UTR}^{(W)}_t$) can be identified.
% This enables mapping of occupancy conditions where Dynamic control is equivalence-effective versus conditions where its advantage is limited.


% \begin{itemize}[noitemsep]

%     \item \textbf{Comfort Coverage (person--time):}
%     For Static, Daily, and Dynamic policies, report absolute coverage (\%) and differences
%     $\Delta\mathrm{Cov}^{\mathrm{Dyn-Static}},\,\Delta\mathrm{Cov}^{\mathrm{Dyn-Daily}}$
%     at room and building levels, and across occupancy bins, identifying the ranges of
%     $N_{\text{avg}}, V_{\text{occ}}, \mathrm{UTR}^{(W)}_t$ where Dynamic yields higher coverage.

%     \item \textbf{Energy Impact (counterfactual):}
%     Holding weather and occupancy fixed, estimate policy-specific total energy and saving rates
%     (Dynamic vs.\ Static/Daily) from \acs{bms}-derived component models.
%     Report totals, peak-related indicators (if applicable), and uncertainty intervals.

%     \item \textbf{Setpoint Equivalence (TOST, $\pm 0.5^\circ$C):}
%     For each room and overall, determine under which occupancy conditions the Dynamic setpoint is statistically equivalent to, or exceeds, Static and Daily baselines. Results are summarized by ordered occupancy bins (e.g., by $N_t$, $\sigma^{(W)}_{N,t}$, $\mathrm{UTR}^{(W)}_t$).

%     \item \textbf{Actionable Decision Rules:}
%     Occupancy-informed rules of thumb---thresholds on $N_{\text{avg}}, V_{\text{occ}}, \mathrm{UTR}^{(W)}_t$---to choose update frequency (Static / Daily / Dynamic) for practical deployment.
% \end{itemize}



% Metrics outcome

% Among the three \acp{gcm} defined above, the \ac{rtgcm} setpoint $T_{\mathrm{set}}^{\mathrm{RT}}(t)$ is treated as the highest-resolution target and is compared with the \ac{sgcm} and \ac{dgcm} baselines.
% The deviations are
% \begin{equation}
% \Delta T^{\mathrm{Static}}_t
% \;=\;
% T_{\mathrm{set}}^{\mathrm{RT}}(t)
% -
% T_{\mathrm{set}}^{\mathrm{Static}}(t),
% \qquad
% \Delta T^{\mathrm{Daily}}_t
% \;=\;
% T_{\mathrm{set}}^{\mathrm{RT}}(t)
% -
% T_{\mathrm{set}}^{\mathrm{Daily}}(t).
% \label{eq:deltaT_defs}
% \end{equation}
% Unless stated otherwise, the magnitudes of $|\Delta T^{\mathrm{Static}}_t|$ and $|\Delta T^{\mathrm{Daily}}_t|$ will be analyzed.

% \begin{figure}[pos=!bt]
%     \centering
%     \includegraphics[width=0.97\linewidth]{\Pic{png}{ex1_analysis_flow}}
%     \caption{Analysis of Baseline vs Dynamic comfort model Image}
%     \label{fig:ex1analysis}
% \end{figure}

\FloatBarrier

\section{Results}
\subsection{Comfort performance across different resolution of \ac{gcm}}

\SetPicSubDir{ch4-Results/agentic/Result_Agentic}
\label{subsec:ComfPerf}

Figure~\ref{fig:ComfAgentic}\subref{fig:ComfAchi} shows the mean comfort probability under the \ac{sgcm}, \ac{dgcm}, \ac{rtgcm}, and fixed \qty{24}{\celsius} baseline policies for different subgroup sizes across the four office rooms.
Figure~\ref{fig:ComfAgentic}\subref{fig:ComfImp} shows the improvement in mean comfort probability relative to the fixed \qty{24}{\celsius} baseline.
Across all rooms, the baseline control remained nearly constant at approximately \qtyrange{65}{70}{\percent} across subgroup sizes.
The \ac{rtgcm} produced the highest mean comfort probability for every room and subgroup size, followed by the \ac{dgcm} and then the \ac{sgcm}.
The performance of the \ac{rtgcm} and \ac{dgcm} decreases at similar rates as subgroup size increases, while the performance of the \ac{sgcm} shows a more modest decrease.
The mean comfort probability across rooms ranged approximately \qtyrange{80}{98}{\percent} for the \ac{rtgcm}, \qtyrange{75}{95}{\percent} for the \ac{dgcm}, and \qtyrange{65}{73}{\percent} for the \ac{sgcm}.

The improvement relative to the 24 \si{\celsius} baseline decreased as subgroup size increased.
\ac{sgcm} improvement approaches zero in larger rooms, especially around subgroup size of 5--7 users.

In contrast, the \ac{rtgcm} and \ac{dgcm} maintain positive improvements at the largest tested subgroup size, with \ac{rtgcm} improvement remaining approximately \qtyrange{15}{20}{\percent} and \ac{dgcm} improvement approximately \qtyrange{5}{20}{\percent}.

Although both the \ac{dgcm} and \ac{rtgcm} maintain performance improvements over the baseline, the gap between them increases as room size expands.
This suggests that, in larger rooms, the \ac{rtgcm} benefits from more dynamic occupant transitions, whereas the \ac{dgcm} cannot fully capture time-varying group comfort as occupancy changes.

\begin{figure}[pos=htbp!]
    \centering
    \begin{subfigure}{0.98\linewidth}
        \centering
        \includegraphics[width=\linewidth]{\Pic{png}{a6p_static_realtime_daily_by_subgroup}}
        \caption{Mean comfort probability under different comfort-representation controls}
        \label{fig:ComfAchi}
    \end{subfigure}

    \medskip

    \begin{subfigure}{0.9\linewidth}
        \centering
        \includegraphics[width=\linewidth]{\Pic{png}{a6p_improvement_vs24_by_subgroup}}
        \caption{Comfort improvement relative to the baseline control}
        \label{fig:ComfImp}
    \end{subfigure}
    \caption{Agentic simulation results for different subgroup sizes.}
    \label{fig:ComfAgentic}
\end{figure}

% \autoref{fig:MeanOccuConf} shows the improvement ratio relative to the baseline control ($24^\circ\mathrm{C}$) as a function of daily mean occupancy. The figure presents the relationship between daily mean occupancy and improvement for the real-time and static group comfort models using density contours and fitted linear regression lines. Most observations are concentrated in regions where the improvement ratio is below 20\%. The fitted lines further indicate that the real-time group comfort model is more effective under sparse occupancy conditions. As mean occupancy increases to around 4--5 occupants, the improvement ratio approaches zero. This suggests that, even when room size (subgroup size) is fixed, the actual mean occupancy level still affects the achievable improvement.

% \begin{figure}[pos=htbp!]
%     \centering
%     \includegraphics[width=0.97\linewidth]{\Pic{png}{analysis72_mean_occupancy_contour_all_k4_k6_k8}}
%     \caption{Improvement ratio relative to the baseline control across daily mean occupancy levels}
%     \label{fig:MeanOccuConf}
% \end{figure}

\subsection{Effect of \ac{rtgcm} on control adjustment}
\label{subsec:controleffect}
The previous subsection showed that the \ac{rtgcm} maintained a \qtyrange{15}{30}{\percent} higher mean comfort probability than the \ac{sgcm} across room member sizes.
This subsection investigates the control-adjustment burden introduced by the \ac{rtgcm}.
Because the \ac{sgcm} and \ac{dgcm} use coarser occupant-composition representations, frequent setpoint updates are primarily required under \ac{rtgcm}-based control.

\autoref{fig:DynSpStepHist} shows the distribution of temperature adjustment magnitudes under the \ac{rtgcm} across daily mean occupancy.
As daily occupancy increases, the range of temperature adjustment magnitude slowly shrinks.
However, the mean value remains approximately around \qty{1.0}{\celsius} even at the highest occupancy levels.

Similarly, \autoref{fig:DynSpStepHeatmap} maps the temperature adjustment magnitude over subgroup size and mean occupancy ratio at the control timestep.
The magnitude of this temperature range shrinks as both subgroup size and mean occupancy ratio increase.
At small subgroup sizes such as $K=4$, the mean adjustment magnitude is over \qty{1.0}{\celsius} in most cases.


\begin{figure}[pos=htbp!]
    \centering
    \includegraphics[width=0.8\linewidth]{\Pic{png}{analysis23_mean_occupancy_moving_quantile_dynamic_sp_mean_abs_step_when_changed_expanded_pool_overlap_dense_only}}
    \caption{Distribution of temperature adjustment magnitudes under dynamic group control}
    \label{fig:DynSpStepHist}
\end{figure}

\begin{figure}[pos=htbp!]
    \centering
    \includegraphics[width=0.8\linewidth]{\Pic{png}{analysis23_k_normalized_mean_occupancy_mean_abs_step_when_changed_heatmap_regime_style_expanded_pool_overlap_dense_ybins05pct_kle16_minn20}}
    \caption{Mapping of temperature adjustment magnitude over subgroup size and mean occupancy ratio}
    \label{fig:DynSpStepHeatmap}
\end{figure}

\subsection{Effect of occupancy dynamics on control action frequency}

\autoref{fig:CtrlWindowRatio} shows the probability of temperature adjustment events of at least \qty{0.5}{\celsius} across different \ac{utr} levels.
As \ac{utr} increases, the probability of required temperature adjustment also increases; when \ac{utr} reaches approximately 0.7--0.8, the probability of temperature adjustment reaches \qty{90}{\percent}.

\autoref{fig:UTRexamples} shows the mean \ac{utr} transition by time of day, calculated as the weekday mean in four offices.
Generally, \ac{utr} marks the highest values around \DTMtime{08:00:00} to \DTMtime{08:30:00} and around \DTMtime{17:30:00} to \DTMtime{18:00:00}, which are typical occupant arrival and departure times, respectively.
Also, \ac{utr} marks higher values at lunchtime around \DTMtime{12:00:00} to \DTMtime{13:00:00} in shared offices and R\&D Office.
In Shared Office 1 and Shared Office 2, the values change more dynamically than in Admin Office C during office hours.
\ac{utr} marks a higher value around 1.0 on Mondays in R\&D Office, and in Shared Office 2 \ac{utr} marks approximately 0.8 on Mondays.
Except these cases, the \ac{utr} values are generally between 0.2 to 0.5.

\begin{figure}[pos=htbp!]
    \centering
    \includegraphics[width=0.9\linewidth]{\Pic{png}{analysis25_utr30_endpoint_k_scaling_step_rate_full_confirm_rng_k16_trial1_1000_combined_20260416_073500}}
    \caption{Relationship between 30-minute user turnover rate and action-needed window ratio}
    \label{fig:CtrlWindowRatio}
\end{figure}

\begin{figure}[pos=htbp!]
    \centering
    \includegraphics[width=0.9\linewidth]{\Pic{png}{analysis25_utr30_rolling_representative_k_separate_room_panels_30min_full_confirm_rng_k16_trial1_1000_combined_20260416_073500}}
    \caption{Distribution of temperature adjustment magnitudes under dynamic group control}
    \label{fig:UTRexamples}
\end{figure}



% removed
% \SetPicSubDir{ch4-Results/agentic/Result_GEARsimulation}
% As a first part of agentic simulation, realistic one-month simulation was conducted.
% This case is based on actual combination of occupancy logs and comfort models with survey participants.
% An example of daily temperature control and predicted comfort performance both on \ac{sgcm} and \ac{rtgcm} is shown in \autoref{fig:OpeExampls}.

% \autoref{fig:GEARComfPrf}\subref{fig:realComfPro} shows the simulated comfort probability on \ac{sgcm} and \ac{rtgcm} during one-month simulation period. \autoref{fig:GEARComfPrf}\subref{fig:realComfImp} shows the predicted comfort probability improvement from \ac{sgcm} to \ac{rtgcm} during one-month simulation period.

% \begin{figure}[pos=htbp!]
%     \centering
%     \begin{subfigure}{0.97\linewidth}
%         \centering
%         \includegraphics[width=\linewidth]{\Pic{png}{ComfortProbability}}
%         \caption{Simulated comfort probability on \ac{sgcm} and \ac{rtgcm} during one-month simulation period}
%         \label{fig:realComfPro}
%     \end{subfigure}

%     \medskip

%     \begin{subfigure}{0.97\linewidth}
%         \centering
%         \includegraphics[width=\linewidth]{\Pic{png}{ComfortProbabilityImprovement}}
%         \caption{Predicted comfort probability improvement from \ac{sgcm} to \ac{rtgcm} during one-month simulation period}
%         \label{fig:realComfImp}
%     \end{subfigure}
%     \caption{Simulated mean comfort probability of each experiment participant through one-month simulation}
%     \label{fig:GEARComfPrf}
% \end{figure}

% \begin{figure}[pos=htbp!]
%     \centering
%     \includegraphics[width=0.97\linewidth]{\Pic{png}{ComfortAcrossRooms}}
%     \caption{Comfort probabiity ratio across office rooms and whole building on \ac{sgcm} and \ac{rtgcm}}
%     \label{fig:GEARComfPrfRoom}
% \end{figure}

\FloatBarrier

\section{Discussion}
\SetPicSubDir{ch5-Discussion}

\subsection{Why aggregation resolution changes the group-size comfort performance}
% \subsection{Effect of occupancy parameters on comfort performance}
\mycite{jung_energy_2020}\cite{jung_energy_2020} investigated how the number of occupants in a thermal zone affects the performance of member-oriented comfort models.
Their study showed that the comfort benefit of personal-comfort-model integration tends to diminish as more occupants share the same zone, with the improvement converging to nearly two percent when the number of occupants increases to six.
At larger group scales, increasing group size and internal preference diversity likewise reduced the potential maximum satisfaction rate, while larger groups became less thermally sensitive~\cite{wang_enhancing_2026}.

This tendency is consistent with the present simulation results.
As shown in \autoref{fig:MeanOccuConf}, the improvement of the \ac{sgcm} relative to the fixed \qty{24}{\celsius} baseline decreases as subgroup size increases and approaches zero around subgroup sizes of five to seven.

These results are consistent with prior multi-occupant comfort-control studies showing that the benefit of integrating personal comfort information tends to diminish as more occupants are aggregated within a single thermal zone.
As the realized group becomes larger, individual thermal preferences are averaged within the group representation, reducing the difference among alternative aggregation resolutions.
At the same time, the \ac{dgcm} and \ac{rtgcm} maintain positive improvement over the fixed \qty{24}{\celsius} baseline even at larger subgroup sizes.
This indicates that the value of higher-resolution \acp{gcm} depends not only on room-member size, but also on realized occupancy and occupant-composition dynamics.

The SGCM here follows the same directional pattern, crossing zero at $K\approx6$--8.
DGCM and RT-GCM, however, retained positive improvement at the same $K$ by excluding absent members.
This study therefore extends prior group-size findings by showing that comfort-aggregation outcomes depend not only on potential group size but also on realized occupancy (Figure~\ref{fig:OccuDensity}) and dynamic changes in occupant composition.
RT-GCM also exceeded DGCM by \qtyrange{6}{8}{\percent} points after the initial rise with $K$, reflecting the value of knowing who is present now rather than who attended that day.
These predicted gains assume immediate setpoint selection and remain conditional on sensing reliability, HVAC response, and occupant adaptation.

\subsection{From comfort information to actionable control}
% \subsection{Effect of occupancy parameters on control}
As examined in \autoref{subsec:controleffect}, the required setpoint adjustment under \ac{rtgcm}-based control can remain operationally distinguishable across a wide range of subgroup sizes and occupancy conditions.
Even under larger subgroup sizes and higher mean occupancy, the adjustment magnitude often exceeds the common thermostat control step of \qty{0.5}{\celsius}.
In this simulation, the control temperature was assumed to be adjustable at a resolution of \qty{0.1}{\celsius}, allowing the \ac{rtgcm} to express finer comfort-driven setpoint changes.
A high-resolution comfort representation is useful only when the HVAC system can respond to the resulting setpoint differences.
Although the simulation used a \qty{0.1}{\celsius} grid, typical changed steps remained near \qty{0.8}{\celsius} at higher mean occupancy and exceeded the \qty{0.5}{\celsius} actionable threshold.
This indicates actuation-relevant potential; reviews report limited field implementation and persistent data, computation, BAS-integration, and scalability barriers field examples~\cite{xie_review_2020,soleimanijavid_challenges_2024}.

This finding also relates to the resolution-mismatch argument reported by \mycite{ono_effects_2022}\cite{ono_effects_2022}.
A high-resolution comfort representation can only be useful in practice when the available \ac{hvac} control resolution can respond to the corresponding setpoint differences.
The present results show this trade-off from the operational side: the \ac{rtgcm} can improve mean comfort probability, but it also introduces larger and more frequent setpoint adjustments.

As examined in \autoref{fig:CtrlWindowRatio}, \ac{utr} has a strong effect on the probability of required temperature adjustment.
Mean occupancy and occupancy variability describe how many people are present and how strongly the count fluctuates, whereas \ac{utr} captures who is replaced over short time intervals.
This distinction is important for \ac{gcm}-based control because the preferred group setpoint can change even when the occupancy count remains similar, if the occupant composition changes.

\autoref{fig:UTRexamples} further shows different turnover patterns between shared offices and seat-assigned offices such as Admin Office C.
Shared offices show more dynamic \ac{utr} transitions during office hours, while assigned offices generally show lower values such as \qtyrange{0.2}{0.4}{}.
Some offices also show weekly turnover patterns; for example, Admin Office C and Shared Office 2 are affected by group meetings outside the target office rooms on Monday afternoon.
\ac{utr} can capture such office-use patterns, indicating its usefulness for identifying when \ac{rtgcm}-based control is likely to be operationally relevant.

\ac{utr} complements mean occupancy by indicating whether represented comfort profiles are being replaced, extending prior discussions of dynamic diversity in comfort-driven control \cite{jung_energy_2020}.
The near-overlap across $K$ in Figure~\ref{fig:CtrlWindowRatio} shows that turnover is a better guide to update frequency than $K$ alone.
The three aggregation levels can consequently be treated as a deployment hierarchy: \ac{sgcm} is defensible in stable assigned-seat zones, where \ac{utr} remained near 0.2--0.4 through office hours and the regular-member group matched presence; \ac{dgcm} is a lower-complexity compromise when daily attendance varies but within-day composition stays relatively stable.
\ac{rtgcm} is most valuable where \ac{utr} instead fluctuated dynamically, with sparse attendance and frequent movement.
\ac{utr} can gate real-time updates to high-turnover periods, with deadbands or minimum hold times limiting excessive cycling.
Such gating does not remove sensing trade-offs in accuracy, latency, energy use, cost, scalability, and privacy~\cite{zafari_survey_2019}.

These findings suggest that \ac{rtgcm}-level control should not be adopted solely because real-time tracking is available.
It is most valuable in sparse or highly variable shared-office conditions where occupant composition changes frequently.
In contrast, a \ac{dgcm} may provide a practical compromise when daily attendance differs from regular membership but within-day turnover is moderate, while a \ac{sgcm} may be sufficient when occupancy is stable or group size is large enough for individual differences to average out.

Because activity-based working outcomes depend on implementation and occupant use~\cite{marzban_review_2023}, these policy thresholds should be calibrated to each workplace.

\subsection{Limitations}
Several limitations should be noted.
First, the data come from one office building and one observation period, so the numerical thresholds for subgroup size, mean occupancy, and \ac{utr} may depend on building use, occupant population, room size, and survey protocol.
Second, the agentic simulation reassigns \ac{pcm} profiles to observed location-log agents.
This increases scenario coverage but does not represent the actual person-specific coupling between movement patterns and comfort preferences.
Third, comfort evaluation is based on predicted ``No change'' probability rather than direct closed-loop field operation.
Energy consumption, \ac{hvac} response delay, actuator constraints, and occupant adaptive behavior were not fully evaluated.
Finally, real-time tracking availability does not automatically imply deployability; privacy, sensing reliability, and integration cost remain practical constraints.

\FloatBarrier

This study investigated how the temporal resolution of \ac{gcm} aggregation affects predicted comfort and control burden in dynamically occupied shared offices.
Three \acp{gcm} were compared: a \ac{sgcm} based on regular room members, a \ac{dgcm} based on daily attendees, and a \ac{rtgcm} based on occupants present at each control decision window.
Using field survey data and reconstructed location logs, the agentic simulation showed that higher-resolution \acp{gcm} can improve mean comfort probability relative to a fixed \qty{24}{\celsius} baseline and the \ac{sgcm}.
\ac{rtgcm} achieved the highest predicted mean comfort probability across all tested rooms and group sizes, decreasing from \qtyrange{92}{93}{\percent} at $K=2$ to \qtyrange{78}{82}{\percent} at the largest tested $K$.
At the largest group sizes, \ac{rtgcm} remained \qtyrange{6}{8}{\percent} gains above \ac{dgcm}, whereas \ac{sgcm} decreased to \qtyrange{64}{66}{\percent}, approximately converging with the \qtyrange{65}{70}{\percent} fixed baseline.
The advantage of the \ac{rtgcm} was strongest under sparse or dynamically changing occupancy conditions.

At the same time, the \ac{rtgcm} introduced operational demands, reflected in setpoint adjustment magnitude and the action-needed window ratio.
The typical \ac{rtgcm} setpoint change decreased from approximately 1.5$^\circ$C at a daily mean occupancy near one to \qty{0.8}{\celsius} at the highest observed occupancy, remaining above the \qty{0.5}{\celsius} actionable threshold.
The actionable-update rate increased from approximately 37\% at a \ac{utr} near 0.05 to \qtyrange{89}{93}{\percent} at a \ac{utr} of 0.9--1.0, with little dependence on $K$.

These findings indicate that high-resolution occupancy tracking should be deployed selectively, based on occupancy metrics such as mean occupancy, subgroup size, occupancy variability, and \ac{utr}.
Future work should include closed-loop field tests, energy analysis, \ac{hvac} dynamics, privacy-aware sensing implementation, and longer deployments across different office types.

Within this comfort-only simulation scope, these results support static aggregation for stable membership, daily aggregation for variable attendance with moderate turnover, and real-time aggregation for sparse, rapidly changing occupancy.
\ac{utr} could therefore serve as a screening metric for selecting the aggregation resolution.
\FloatBarrier

% \begin{figure}[pos=hbt!]
%     \centering
%     \includegraphics[width=0.95\linewidth]{figure/3.Method/EMLdiagram.png}
%     \caption{Data input, calculation and Evaluation flow}
%     % todo - FT - could you please add more information to the caption?->I refined but am not confident if it follows what you expect
%     \label{fig:AnaFlow}
% \end{figure}




\FloatBarrier

\subsection{Limitations}
XXX

\acresetall
XXX

% \section{Cross-references}
% In electronic publications, articles may be internally
% hyperlinked. Hyperlinks are generated from proper
% cross-references in the article.  For example, the words
% \textcolor{black!80}{Fig.~1} will never be more than simple text,
% whereas the proper cross-reference \verb+\ref{tiger}+ may be
% turned into a hyperlink to the figure itself:
% \textcolor{blue}{Fig.~1}.  In the same way,
% the words \textcolor{blue}{Ref.~[1]} will fail to turn into a
% hyperlink; the proper cross-reference is \verb+\cite{Knuth96}+.
% Cross-referencing is possible in \LaTeX{} for sections,
% subsections, formulae, figures, tables, and literature
% references.

% \section{Bibliography}

% Two bibliographic style files (\verb+*.bst+) are provided ---
% {model1-num-names.bst} and {model2-names.bst} --- the first one can be
% used for the numbered scheme. This can also be used for the numbered
% with new options of {natbib.sty}. The second one is for the author year
% scheme. When  you use model2-names.bst, the citation commands will be
% like \verb+\citep+,  \verb+\citet+, \verb+\citealt+ etc. However when
% you use model1-num-names.bst, you may use only \verb+\cite+ command.

% \verb+thebibliography+ environment.  Each reference is a\linebreak
% \verb+\bibitem+ and each \verb+\bibitem+ is identified by a label,
% by which it can be cited in the text:

% \noindent In connection with cross-referencing and
% possible future hyperlinking it is not a good idea to collect
% more that one literature item in one \verb+\bibitem+.  The
% so-called Harvard or author-year style of referencing is enabled
% by the \LaTeX{} package {natbib}. With this package the
% literature can be cited as follows:

% \begin{enumerate}[\textbullet]
% \item Parenthetical: \verb+\citep{WB96}+ produces (Wettig \& Brown, 1996).
% \item Textual: \verb+\citet{ESG96}+ produces Elson et al. (1996).
% \item An affix and part of a reference:\break
% \verb+\citep[e.g.][Ch. 2]{Gea97}+ produces (e.g. Governato et
% al., 1997, Ch. 2).
% \end{enumerate}

% In the numbered scheme of citation, \verb+\cite{<label>}+ is used,
% since \verb+\citep+ or \verb+\citet+ has no relevance in the numbered
% scheme.  {natbib} package is loaded by {cas-sc} with
% \verb+numbers+ as default option.  You can change this to author-year
% or harvard scheme by adding option \verb+authoryear+ in the class
% loading command.  If you want to use more options of the {natbib}
% package, you can do so with the \verb+\biboptions+ command.  For
% details of various options of the {natbib} package, please take a
% look at the {natbib} documentation, which is part of any standard
% \LaTeX{} installation.

\pagebreak

% Loading bibliography database
\bibliography{references}

\pagebreak

\appendix
\label{appendix:XXX}

\SetPicSubDir{chap-appe}

\subsection{Office tracking and visualization technology adoption cases}
\label{app:Tracking}
\autoref{table:OccuTrack} shows examples of offices that utilize worker's location information during operation.

% \begin{table}[pos=hbt!]
%   \centering
%   \centering
\small
\setlength{\tabcolsep}{4pt}
\renewcommand{\arraystretch}{1.2}
\begin{tabularx}{\textwidth}{l p{3cm} p{4cm} p{6cm}}
\toprule
\textbf{Ref.} & \textbf{Company} & \textbf{Office} & \textbf{Technology} \\
\midrule
\cite{nihonsekkeiinc.RealtimeVisualization2024} & Nihon Sekkei & Head Office (Tokyo) & Smartphone + BLE beacons (Digital Twin visualization) \\
\cite{the_gear_kajima_2024} & Kajima Corporation & Singapore HQ ``The GEAR'' & IoT sensors + AI analytics\\
\cite{itokiITOKI2011} & Itoki & ITOKI TOKYO XORK & BLE beacons + mobile app \\
\cite{kokuyoco.OneYearRecords2021} & Kokuyo & Head Office & BLE beacons + mobile app \\
\cite{uchidayokoPressReleaseUchidaYoko2024} & Uchida Yoko & Shinkawa HQ Office, Olympus HQ, etc. & Wi-Fi indoor positioning + office map visualization \\
%\cite{colorkrew_biz_colorkrew_nodate} & Marubeni Corp. & HQ Chemical Division & QR code desk management (Colorkrew Biz) \\
%\cite{acall_it_2022} & Hashimoto Gumi & HQ Building & QR code desk management (Acall) \\
\cite{tanaka_development_2022} & Takenaka Corporation & Takenaka Central Bldg. South & AI cameras (human distribution \& posture sensing) \\
\bottomrule
\end{tabularx}

%   \caption{Office tracking and visualization technology adoption cases in Japan and Singapore}
%   \label{table:OccuTrack}
% \end{table}

\subsection{Occupancy Log Reconstruction Procedure}
\label{app:occupancy_reconstruction}

The raw data stream for occupancy tracking contains three event types from a face-recognition system:
\begin{itemize}[noitemsep]
    \item \textbf{Room Entry Events:} Triggered when a participant enters a specific room through a door-mounted access point.
    \item \textbf{General Presence Events:} Triggered when a participant is detected in shared zones or corridors, confirming their presence in the building.
    \item \textbf{Building Exit Events:} Triggered when a participant leaves the building through a main egress gate. Note that room-level exit events are not always recorded.
\end{itemize}

From these discrete events, a structured ``stay log'' is generated to represent continuous periods of presence. For each participant and each day, the following rules are applied:
\begin{itemize}[noitemsep]
    \item A stay begins at the timestamp of an “Entry” or equivalent presence-confirming event, and ends at the next recorded event (room entry, zone presence, or building exit).
    \item To prevent overnight carryover, any stay still active at the day’s final event is terminated at that time, typically corresponding to a building exit.
\end{itemize}

This process yields a stay log with anonymized IDs, locations, and start–end times for subsequent analysis.

\subsection{Occupancy tracking studies and examined parameters}
\label{app:OccupancyStudy}
Below \autoref{table:OccuMetrics} list up research studies that examined dynamic occupancy and the parameters used to measure dynamics.

\begin{table}[pos=htbp!]
  \centering
  \input{pic/chap-appe/ExOccuStudy}
  \caption{Summary of occupancy-related literature and metrics}
  \label{table:OccuMetrics}
\end{table}

\subsection{Thermal comfort survey contents}
\label{app:ThermalSurvey}
Below \autoref{table:comfsurvey} list up research studies that examined dynamic occupancy and the parameters used to measure dynamics.

\begin{table}[pos=htbp!]
  \centering
  \input{pic/chap-appe/ThermalComfort}
  \caption{Contents of Thermal comfort survey}
  \label{table:comfsurvey}
\end{table}


\SetPicSubDir{ch4-Results}
\subsection{Thermal comfort survey contents}
\begin{figure}[pos=htbp!]
    \centering
    \includegraphics[width=1.0\linewidth]{\Pic{png}{/agentic/Result_GEARsimulation/OperationExample}}
    \caption{Example of daily temperature control and predicted comfort performance both on \ac{sgcm} and \ac{rtgcm}}
    \label{fig:OpeExampls}
\end{figure}
\FloatBarrier


% Appendix sections are coded under \verb+\appendix+.

% \verb+\printcredits+ command is used after appendix sections to list 
% author credit taxonomy contribution roles tagged using \verb+\credit+ 
% in frontmatter.

\setcounter{table}{0}
\setcounter{figure}{0}
\renewcommand{\thetable}{A.\arabic{table}}
\renewcommand{\thefigure}{A.\arabic{figure}}

XXX

\section{XXX}
\label{appendix:XXX}

\renewcommand{\thetable}{B.\arabic{table}}
\renewcommand{\thefigure}{B.\arabic{figure}}


%\vskip3pt

% \bio{}
% Author biography without author photo.
% Author biography. Author biography. Author biography.
% \endbio

% \bio{figs/pic1}
% Author biography with author photo.
% Author biography. Author biography. Author biography.
% \endbio

% \bio{figs/pic1}
% Author biography with author photo.
% Author biography. Author biography. Author biography.
% \endbio

\end{document}
